\documentclass[conference]{IEEEtran}
\IEEEoverridecommandlockouts

\usepackage{cite}
\usepackage{amsmath,amssymb,amsfonts}
\usepackage{algorithmic}
\usepackage{graphicx}
\usepackage{textcomp}
\usepackage{xcolor}
\usepackage{booktabs}
\usepackage{multirow}
\usepackage{array}

\def\BibTeX{{\rm B\kern-.05em{\sc i\kern-.025em b}\kern-.08em
    T\kern-.1667em\lower.7ex\hbox{E}\kern-.125emX}}

\begin{document}

\title{From Learning Interactions to Pedagogically Structured Evidence:\\A Distributional Framework for Knowledge Tracing}

\author{%
\IEEEauthorblockN{Author One}
\IEEEauthorblockA{\textit{Department of Computer Science} \\
\textit{University Name}\\
City, Country \\
author1@university.edu}
\and
\IEEEauthorblockN{Author Two}
\IEEEauthorblockA{\textit{Department of Computer Science} \\
\textit{University Name}\\
City, Country \\
author2@university.edu}
}

\maketitle

\begin{abstract}
Knowledge Tracing (KT) aims to model the evolution of student knowledge states from historical interaction sequences to predict future performance. While deep KT models have achieved strong predictive accuracy, they typically update mastery representations directly from interaction encodings, bypassing the structured pedagogical evidence that learning interactions inherently contain. In this paper, we propose a distributional framework that explicitly introduces an intermediate stage between interaction encoding and mastery updating. Our framework consists of four modules: (1)~Interaction Representation Learning, which captures dual dynamics of interaction sequences and localized knowledge evolution via parallel encoders; (2)~Distributional Pedagogical Alignment, which projects interaction representations into a probabilistic space and constructs multi-view learning patterns with pedagogical constraints; (3)~Mastery State Tracking, which updates knowledge states from structured patterns while quantifying pattern-to-knowledge contributions; and (4)~Prediction Aggregation, which generates predictions with knowledge-component attribution for intrinsic interpretability. The key insight is that pedagogical evidence should be organized, aligned, and structured before informing mastery updates, rather than being implicitly compressed into latent representations. This design provides a principled separation of representation spaces and enables traceable reasoning from interaction patterns to knowledge states to predictions. We present the framework design and discuss its theoretical foundations and implications for interpretable knowledge tracing.
\end{abstract}

\begin{IEEEkeywords}
Knowledge Tracing, Distributional Representation, Learning Pattern, Pedagogical Alignment, Interpretable AI
\end{IEEEkeywords}

% ====================================================================
\section{Introduction}
\label{sec:intro}

\subsection{Background}

Knowledge Tracing (KT) addresses the fundamental problem of modeling how student knowledge evolves through a sequence of learning interactions, with the goal of predicting performance on future tasks~\cite{corbett1995bkt, piech2015dkt}. Over the past decade, deep learning-based KT models have achieved remarkable predictive performance by leveraging advanced sequence modeling architectures. Deep Knowledge Tracing (DKT)~\cite{piech2015dkt} pioneered the use of recurrent neural networks for this task, while subsequent works have introduced attention mechanisms~\cite{pandey2019akt}, graph structures~\cite{nakagawa2019gkt}, and state-space models~\cite{li2025mckt} to capture richer interaction dynamics.

More recently,KT research has expanded along three complementary directions: (i)~enhancing interaction representation through content-aware and topology-aware encoding~\cite{li2025mckt, xia2025}; (ii)~modeling knowledge states as distributions rather than deterministic vectors to capture uncertainty~\cite{shen2024sskt, liu2025keenkt, chen2025ukt, yoon2018deepirt}; and (iii)~improving interpretability through neural-symbolic reasoning~\cite{zhang2023nskt} or probabilistic inference~\cite{ghosh2020psikt}. Additionally, temporal dynamics modeling has been explored through forgetting mechanisms~\cite{zhang2019dktforget} and point process-based approaches~\cite{li2020hawkeskt}, while hierarchical attention architectures~\cite{pandey2023dtransformer} and pedagogical theory-driven models~\cite{cui2024grkt} have emerged as promising directions. Despite these advances, most existing models share a common architectural assumption: the interaction representation learned by the encoder is used \emph{directly} to update the knowledge state.

\subsection{Research Gap}

We identify three interconnected limitations in this direct-update paradigm.

\textbf{Gap 1: Absence of explicit pedagogical evidence.} Existing KT models compress all relevant information---interaction semantics, difficulty, student response, and historical mastery---into a single latent representation. Pedagogical signals (e.g., whether a student is showing a declining trend, has mastered prerequisites, or is struggling with a specific concept) are implicitly encoded but never explicitly organized or structured.

\textbf{Gap 2: No intermediate reasoning stage.} The mapping from interaction representation to knowledge state update is typically a single neural transformation. There is no intermediate stage where evidence can be aggregated across multiple views (temporal, conceptual, structural), compared, or aligned with pedagogical principles before informing mastery changes.

\textbf{Gap 3: Limited traceability of knowledge updates.} When a knowledge state is updated, it is difficult to determine \emph{which} learning patterns contributed to the change, or \emph{how much} each pattern influenced specific knowledge components. This makes the knowledge updating process essentially opaque.

These limitations are illustrated in Figure~\ref{fig:motivation}.

\begin{figure}[htbp]
\centering
\small
\begin{tabular}{c@{\hspace{4pt}}c}
\textbf{Existing} & \textbf{Ours} \\
\begin{minipage}{0.35\columnwidth}
\centering
\begin{tabular}{c}
\textit{Interaction} \\
$\downarrow$ \\
\textit{Mastery} \\
$\downarrow$ \\
\textit{Prediction}
\end{tabular}
\end{minipage}
&
\begin{minipage}{0.55\columnwidth}
\centering
\begin{tabular}{c}
\textit{Interaction} \\
$\downarrow$ \\
\textit{Pedagogical Evidence} \\
$\downarrow$ \\
\textit{Knowledge} \\
$\downarrow$ \\
\textit{Prediction}
\end{tabular}
\end{minipage}
\end{tabular}
\caption{Comparison of existing direct-update paradigm (left) versus our framework with an intermediate pedagogical reasoning stage (right).}
\label{fig:motivation}
\end{figure}

\subsection{Research Hypothesis}

Our central hypothesis is: \emph{Student knowledge should be updated from structured pedagogical evidence---organized, aligned, and quantified learning patterns---rather than directly from isolated interaction representations.}

\subsection{Contributions}

This paper makes the following contributions:

\begin{itemize}
    \item We propose a four-module framework that explicitly separates interaction representation, pedagogical evidence construction, knowledge state tracking, and prediction into distinct representation spaces (Section~\ref{sec:framework}).
    \item We introduce \textit{Distributional Pedagogical Alignment}, a novel intermediate stage that projects interaction representations into a probabilistic space for multi-view pattern construction with pedagogical constraints.
    \item We design intrinsic interpretability mechanisms through pattern-to-knowledge contribution and knowledge-to-prediction contribution, enabling traceable reasoning from patterns to predictions.
\end{itemize}

% ====================================================================
\section{Proposed Framework}
\label{sec:framework}

Our framework organizes the knowledge tracing process into four consecutive representation spaces, each with a distinct role (Figure~\ref{fig:framework}).

\begin{figure}[htbp]
\centering
\small
\begin{tabular}{c}
\textit{Historical Interactions} \\
$\downarrow$ \\
\framebox[0.9\columnwidth]{\textbf{Interaction Representation Learning}} \\
\smallskip
$\downarrow$ \\
\framebox[0.9\columnwidth]{\textbf{Distributional Pedagogical Alignment}} \\
\smallskip
$\downarrow$ \\
\framebox[0.9\columnwidth]{\textbf{Mastery State Tracking}} \\
\smallskip
$\downarrow$ \\
\framebox[0.9\columnwidth]{\textbf{Prediction Aggregation}} \\
\smallskip
$\downarrow$ \\
\textit{Response Probability}
\end{tabular}
\caption{Overall framework: four modules corresponding to four representation spaces.}
\label{fig:framework}
\end{figure}

\subsection{Interaction Representation Learning}
\label{sec:irl}

The first module learns contextual representations of each interaction from the historical sequence. Following recent advances in content-aware and topology-aware KT~\cite{li2025mckt, xia2025}, we separate the input into two groups with distinct dynamics:

\begin{itemize}
    \item \textbf{Interaction Context} -- question semantics, student response, and question difficulty, capturing the dynamics of the interaction sequence.
    \item \textbf{Knowledge Context} -- mastery memory, knowledge graph topology, and concept difficulty, capturing the localized evolution of knowledge state.
\end{itemize}

The Knowledge Context does not access the full mastery memory directly. Instead, it first queries the knowledge graph to retrieve mastery information for concepts relevant to the current question, forming a \emph{localized mastery} representation. This prevents knowledge context from being diluted by irrelevant historical mastery.

Two parallel sequence encoders process these groups independently, and their outputs are fused to form the interaction representation:

\begin{equation}
\mathbf{z}_t = \text{Fusion}\big(\text{Encoder}_{\text{obs}}(I_t),\; \text{Encoder}_{\text{know}}(K_t)\big)
\label{eq:irl}
\end{equation}

where $I_t$ and $K_t$ denote the Interaction Context and Knowledge Context at time step $t$, respectively.

\subsection{Distributional Pedagogical Alignment}
\label{sec:dpa}

This module is the core contribution of our framework. It transforms interaction representations into a probabilistic space where learning patterns can be constructed, aggregated, and aligned with pedagogical principles before informing mastery updates.

Unlike prior works that use distributions to model mastery belief~\cite{ghosh2020psikt} or learner uncertainty~\cite{shen2024sskt, liu2025keenkt}, we use distributions as an \emph{intermediate representation space for interaction patterns}. The module consists of four steps.

\subsubsection{Distribution Projection}

Each interaction representation is projected into a distribution parameterized by a Statistical Descriptor Decoder:

\begin{equation}
h_t = \Phi(\mathbf{z}_t), \quad h_t = (\alpha_t, \beta_t)
\label{eq:projection}
\end{equation}

where $\Phi$ maps the vector representation to distribution parameters. We adopt the Beta distribution, which is naturally bounded on $[0,1]$ and directly interpretable as a probability with associated confidence.

\subsubsection{Pattern Construction}

From the sequence of distributions, we construct multiple learning patterns through multi-view pattern operators. Each operator aggregates distributions along a different semantic axis:

\begin{itemize}
    \item \textbf{Temporal view} -- aggregates patterns over time to capture learning trajectories and progression trends, complementing existing temporal modeling approaches~\cite{li2020hawkeskt, zhang2019dktforget}.
    \item \textbf{Concept-based view} -- aggregates patterns for the same knowledge component to track concept-specific mastery evolution.
    \item \textbf{Structural view} -- aggregates patterns across prerequisite and neighboring concepts to model knowledge transfer, extending hierarchical attention mechanisms~\cite{pandey2023dtransformer}.
\end{itemize}

Each operator produces a pattern representation with consistent semantic meaning, ensuring structural uniformity across students regardless of interaction history length:

\begin{equation}
P_i = \mathcal{A}_i(\mathcal{H})
\label{eq:pattern}
\end{equation}

where $\mathcal{A}_i$ denotes the $i$-th pattern operator and $\mathcal{H}$ is the sequence of distribution representations.

\subsubsection{Pedagogical Alignment}

After pattern construction, soft pedagogical rules adjust the distribution representations to align with educational principles. Unlike NSKT~\cite{zhang2023nskt}, where rules directly determine mastery labels, our rules serve as a \emph{regularization mechanism} on the pattern space:

\begin{equation}
\tilde{P}_i = \text{Align}(P_i, R_i)
\label{eq:align}
\end{equation}

where $R_i$ denotes the set of pedagogical rules applied to pattern $i$. This injects prior pedagogical knowledge into the evidence construction process without deterministically controlling mastery outcomes.

\subsubsection{Pattern Readout}

The adjusted patterns are mapped back to vector representations through a Distribution Readout module. Crucially, the readout produces a \emph{structured representation} with distinct semantic blocks corresponding to each pattern view:

\begin{equation}
\mathbf{z}' = [z^{\text{temp}}, z^{\text{same}}, z^{\text{pre}}, z^{\text{neighbor}}]
\label{eq:readout}
\end{equation}

This structured format ensures that downstream mastery updating receives inputs with consistent structure and meaning across all students.

\subsection{Mastery State Tracking}
\label{sec:mst}

After pedagogical alignment, the framework updates the knowledge state from structured pattern representations. This module serves as the \emph{reasoning stage} where learning patterns are synthesized into mastery changes.

The key mechanism is a \textit{Pattern-to-Knowledge Gating Function} that quantifies the contribution of each learning pattern to each knowledge component:

\begin{equation}
A_i = G(P_i)
\label{eq:gating}
\end{equation}

where $A_i$ represents the influence of the $i$-th pattern on related knowledge components. The mastery update is then computed as a gated combination:

\begin{equation}
\Delta M = \sum_i A_i \odot U_i(M_t, P_i)
\label{eq:update}
\end{equation}

and the knowledge state is updated via residual connection:

\begin{equation}
M_{t+1} = M_t + \Delta M
\label{eq:residual}
\end{equation}

The contribution weights $A_i$ do not participate in prediction but serve as \emph{intrinsic interpretability signals}, enabling questions such as: \emph{Which learning pattern increased the mastery of a specific knowledge component?} and \emph{Which knowledge component is most influenced by a given pattern?}

\subsection{Prediction Aggregation}
\label{sec:pred}

The final module generates predictions from the updated knowledge state and the target question representation. Beyond producing a prediction probability, it simultaneously learns the \textit{Knowledge-Component Contribution}:

\begin{equation}
W = C(M_{t+1}, q_{t+1})
\label{eq:kc_contrib}
\end{equation}

where $W = [w_1, w_2, \ldots, w_K]$ represents the contribution of each knowledge component to the target interaction. The prediction is then:

\begin{equation}
\hat{y} = f(M_{t+1}, q_{t+1}, W)
\label{eq:pred}
\end{equation}

The combination of Pattern-to-Knowledge Contribution (from Module~3) and Knowledge-to-Prediction Contribution (from Module~4) enables \emph{full traceability} of the reasoning chain:

\begin{center}
\small
\begin{tabular}{c}
Learning Patterns \\
$\downarrow$ \\
Pattern $\to$ Knowledge Contribution \\
$\downarrow$ \\
Updated Knowledge State \\
$\downarrow$ \\
Knowledge $\to$ Prediction Contribution \\
$\downarrow$ \\
Prediction
\end{tabular}
\end{center}

This traceability is \emph{intrinsic} to the framework---it emerges from the architecture itself rather than being appended as a post-hoc explanation module.

% ====================================================================
\section{Discussion}
\label{sec:discussion}

\subsection{Why Distribution for Pattern Representation?}

Distributional representations have been explored in KT for modeling mastery belief~\cite{ghosh2020psikt} and learner uncertainty~\cite{shen2024sskt, liu2025keenkt}. Our work repurposes distributions for a fundamentally different goal: representing \emph{interaction evidence} rather than \emph{knowledge state}. A distribution carries not only a central tendency (expectation) but also spread and shape information, enabling richer pattern aggregation and comparison operations than point vectors. However, the distributions in our framework are intermediate representations---they facilitate pattern construction but do not directly encode mastery.

\subsection{Why Pedagogical Alignment before Mastery Update?}

Neural-symbolic approaches like NSKT~\cite{zhang2023nskt} and pedagogical theory-driven models like GRKT~\cite{cui2024grkt} demonstrate the value of incorporating pedagogical rules into KT. Our framework builds on this insight but relocates the role of rules: instead of directly determining mastery outcomes, rules regularize the pattern representation space. This separation ensures that pedagogical knowledge guides evidence construction without overriding the data-driven mastery updating process.

\subsection{Why Intrinsic Interpretability?}

Existing interpretable KT models~\cite{ghosh2020psikt, zhang2023nskt, cui2024grkt} typically achieve interpretability through post-hoc analysis, attention visualization, or pedagogical process modeling. Our framework achieves interpretability \emph{by design}: the pattern-to-knowledge and knowledge-to-prediction contribution mechanisms are integral parts of the architecture. This enables systematic traceability from interaction patterns through knowledge states to predictions, without requiring additional explanation modules.

% ====================================================================
\section{Conclusion}
\label{sec:conclusion}

We have presented a distributional framework for knowledge tracing that explicitly introduces an intermediate pedagogical reasoning stage between interaction encoding and mastery updating. By separating the process into four distinct representation spaces---Interaction, Distribution, Knowledge, and Prediction---our framework provides a principled architecture where pedagogical evidence is constructed, aligned, and structured before informing knowledge state updates. The intrinsic interpretability mechanisms, realized through pattern-to-knowledge and knowledge-to-prediction contribution tracking, enable full traceability of the reasoning chain. This work opens several directions for future research, including distributional reasoning over learning patterns, pedagogically grounded knowledge tracing, and interpretable modeling of knowledge evolution dynamics.

% ====================================================================
% References
% ====================================================================
\bibliographystyle{IEEEtran}
\bibliography{references}

\end{document}
